\section{Related Work}
\label{sec:related}

\paragraph{Engagement detection in educational settings.}
Early work by Whitehill \etal~\cite{whitehill2014faces} demonstrated
that facial features encode cognitive states relevant to learning.
The release of DAiSEE~\cite{gupta2016daisee} established the first large
publicly-available benchmark for four-level engagement classification
from student video.  Subsequent supervised methods leveraged spatiotemporal
features from video~\cite{liao2021deep}, recurrent networks over facial
action units~\cite{thomas2018predicting}, and attention-based
architectures~\cite{huang2019fine}.  More recently, spatiotemporal and attention-based
models trained on DAiSEE have reported accuracy of 55--65\%~\cite{liao2021deep,huang2019fine}, underscoring
the difficulty even for task-specific supervised models.
Classroom-level engagement has received comparatively less attention,
though several works address student behaviour recognition from overhead
cameras~\cite{wang2023scb,kaur2018classroom}.

\paragraph{Zero-shot visual recognition with VLMs.}
CLIP~\cite{radford2021learning} pioneered zero-shot image classification
via image-text cosine similarity, demonstrating competitive performance
on many standard benchmarks.  BLIP~\cite{li2022blip} introduced
bootstrapped image-text alignment that enabled visual question answering
without task-specific training.  Instruction-tuned generative models
such as LLaVA~\cite{liu2023visual} and GPT-4V/4o~\cite{openai2024gpt4o}
further extend this to open-ended visual reasoning.  A growing body of
work has applied these models to medical imaging~\cite{zhang2023pmc},
document understanding, and remote sensing, but their application to
affect recognition and educational observation has not been
systematically studied.

\paragraph{Prompt design and sensitivity.}
The sensitivity of VLMs to prompt phrasing has been studied in the
context of natural language processing~\cite{zhao2021calibrate,
lu2022fantastically}, but visual prompting for classification remains
less understood.  Studies on zero-shot classification~\cite{menon2022visual}
demonstrate that performance can vary by 10--30 percentage points across
prompt phrasings---a phenomenon we observe and quantify here.
